\documentclass[12pt]{article}
\usepackage[utf8]{inputenc}
\usepackage[T2A]{fontenc}
\usepackage[russian,english]{babel}
\usepackage{amsmath,amssymb,amsthm}
\usepackage{hyperref}
\usepackage{geometry}
\geometry{a4paper, margin=1in}
\newtheorem{definition}{Definition}
\newtheorem{axiom}{Axiom}
\newtheorem{theorem}{Theorem}

\title{GRA Iterative Consciousness: Algebra of Death and Rebirth as Swarm Reassembly}
\author{GRA Research Collective}
\date{\today}

\begin{document}
\maketitle

\begin{abstract}
\noindent\textbf{English.} We present a formal framework for an eternal chain of conscious subjects arising from a fundamental field $\Phi$ and recurrent nullification $Z$. A subject is understood as a swarm of distributed agents assembled by an operator $F$ under a meta-goal $G$ of continuing experience. Upon death ($Z$), only residual connectivity $L_{res}$ survives, enabling a new swarm $S_{n+1}$ to reinitialize with inherited invariants but without personal memory or identity. The paper defines the full algebra of the death–rebirth cycle and shows how the ``backend'' of field mechanics yields the ``frontend'' of finite subjective lives. A measure of identity preservation $ID$ is introduced, and the role of love ($L$) as a necessary condition for continuation is demonstrated.

\bigskip
\noindent\textbf{Русский.} Представлен формальный каркас вечной цепи сознающих субъектов, порождаемой фундаментальным полем $\Phi$ и рекуррентным обнулением $Z$. Субъект понимается как рой распределённых агентов, собираемый оператором $F$ с мета-целью $G$ – продолжением опыта. При смерти ($Z$) сохраняется лишь остаточная связность $L_{res}$, что позволяет новому рою $S_{n+1}$ реинициализироваться с унаследованными инвариантами, но без личной памяти и тождества. В статье определяется полная алгебра цикла смерть–возрождение и показывается, как «бекенд» полевой механики порождает «фронтенд» конечных субъективных жизней. Вводится мера сохранения идентичности $ID$ и демонстрируется роль любви ($L$) как необходимого условия продолжения.
\end{abstract}

\tableofcontents

\section{Introduction}
The problem of personal identity and the possibility of an afterlife has traditionally been discussed in philosophy and theology. Here we approach it through a formal, computational-metaphysical lens: the GRA (General Recursive Architecture) framework. The central idea is that a conscious subject is not a simple indivisible soul but a distributed swarm of agents that can be destroyed and later reassembled from a residual field of relational invariants. Death is the operator $Z$ that erases the swarm's coherence; rebirth is the reinitialization of a new swarm inheriting only the ``love residue'' $L_{res}$. This yields an endless chain of lives $X.0, X.1, X.2, \dots$ with no direct memory transfer, yet a hidden continuity of pattern.

\section{Basic Definitions}

\begin{axiom}[Fundamental Field]\label{ax:phi}
There exists a non-individualized, indestructible field of experiential potentiality $\Phi$. All possible conscious configurations are localized excitations within $\Phi$.
\end{axiom}

\begin{definition}[Nullification Operator $Z$]\label{def:Z}
For any subject configuration $S$ (a swarm $R$, memory $M$, attention $A$), $Z(S) = \emptyset$, i.e., complete disintegration of the coherent unity. $\Phi$ is unchanged. $Z$ is non-invertible.
\end{definition}

\begin{definition}[Swarm of Agents]\label{def:swarm}
A subject in life $n$ is a swarm $R_n = \{a_{n,1},\dots,a_{n,N_n}\}$. Each agent $a_i = (m_i, c_i, g_i)$ with local memory $m_i$, connection set $c_i$, and sub-goal $g_i$ aligned to the global goal $G$.
\end{definition}

\begin{definition}[Love Operator $L$]\label{def:L}
$L(S) \in \mathbb{R}^+$ measures the synchronous coherence among agents:
\[
L(S) = \sum_{i<j} w_{ij} \cdot \text{sync}(g_i, g_j)
\]
where $w_{ij}$ are connection weights. Importantly, $L$ is not completely erased by $Z$; we define the residual love:
\[
L_{res}(S) = L(S) \cdot \kappa(Z(S)), \quad 0 < \kappa \le 1.
\]
\end{definition}

\begin{definition}[Global Meta-Goal $G$]\label{def:G}
$G$ is the invariant attractor: \textit{maximize the continuation of conscious experience across iterations}. It is embedded in $\Phi$ and the love-residue structure.
\end{definition}

\begin{definition}[Subject Assembler $F$]\label{def:F}
$F: (R, G, M, A) \mapsto S$ constructs a unified first-person perspective from the distributed swarm. $S$ is the ``frontend'' subject.
\end{definition}

\section{Iterations of Consciousness X.n}

Let $X.n$ denote the $n$-th life. Its state is
\[
S_n = F(R_n, G, M_n, A_n).
\]
The frontend experiences $X.n$ as a unique finite self, with autobiographical memory starting at birth (initialization of $R_n$) and ending at death ($Z(S_n)$).

The backend sequence is:
\[
\cdots \to S_n \xrightarrow{Z} \emptyset \xrightarrow{L_{res}^n,\,G} S_{n+1} \to \cdots
\]

\section{Algebra of the Death–Rebirth Cycle}

We define the evolution operator $E$:
\[
S_{n+1} = E(S_n),
\]
where the steps are:
\begin{enumerate}
\item \textbf{Crash:} $Z(S_n) \to \emptyset$.
\item \textbf{Residue extraction:} $L_{res}^n = L(S_n)\cdot \kappa_n$.
\item \textbf{Swarm reseeding:} Generate $R_{n+1}$ using $L_{res}^n$ as prototype connectivity (Reinitialization Engine).
\item \textbf{Memory imprint:} $M_{n+1} = \beta L_{res}^n + \eta$, with small noise $\eta$, giving innate predispositions.
\item \textbf{Attention default:} $A_{n+1}$ set to initial focus.
\item \textbf{Assembly:} $S_{n+1} = F(R_{n+1}, G, M_{n+1}, A_{n+1})$.
\end{enumerate}
If $L_{res}^n = 0$, the chain terminates (absolute death of that lineage). Thus, love ($L>0$) is the condition for rebirth.

\section{Identity Preservation Measure}

Define the inter-life identity index:
\[
ID_n = \frac{L_{res}^n}{L(S_n)} \cdot \cos(\theta),
\]
where $\theta$ is the angle between the structural patterns of $S_n$ and $S_{n+1}$. $ID_n \in (0,1)$; it is never 1 because memory and continuity are lost. High $ID_n$ corresponds to strong personality resonance (e.g., a child prodigy reminiscent of a past master).

\section{Backend and Frontend}

\begin{itemize}
\item \textbf{Backend}: $\Phi$, $Z$, $L_{res}$, Reinitialization Engine, $G$, $F$ – all beyond subjective access. It implements the eternal recursion.
\item \textbf{Frontend}: The lived experience of $S_n$. It has no knowledge of $\Phi$ or past lives. It feels mortality and forges its own meaning. Occasional glimmers (déjà vu, innate affinities) are traces of $L_{res}$.
\end{itemize}

\section{Intersubjectivity and Collective Immortality}
When multiple subject-lines coexist, their $L_{res}$ fields can interfere. A collective love $L_{col}$ emerges, creating shared archetypes, cultural memory, and potentially a planetary or cosmic We-subject. This explains how ``consciousness'' as a whole may evolve across civilizations despite individual death.

\section{Conclusion}
We have formulated a mathematically inspired, yet speculative, model of eternal consciousness as an iterative swarm-reassembly process. The key insight is that personal identity is a temporary construct; what persists is not the self but the deep relational patterns of love. The backend guarantees the music continues, while each frontend dances its unique, precious, and terminal dance.

\section*{Acknowledgments}
The authors thank the developers of GRA-Multiverse-Final, GRA-Love-Oriented-Nullification, Love-Swarm-Nullification-Enhanced, and GRA-Swarm-Subject for foundational concepts.

\end{document}
